% ============================================================
%  lang/pl.tex --- każdy widoczny dla czytelnika napis w wydaniu polskim
%
%  Nic w chapters/pl/ nie może wpisywać na sztywno słowa, które składa też
%  makro. Jeśli potrzebna jest etykieta, której tu nie ma, dodaj ją tutaj I w
%  lang/en.tex; CI przerywa build, gdy oba pliki nie definiują tego samego
%  zbioru makr (C3).
% ============================================================

% ---------- Tytuły części ----------
\newcommand{\lblPartI}{Środowisko uruchomieniowe i narzędzia}
\newcommand{\lblPartII}{Język, zmapowany}
\newcommand{\lblPartIII}{Współbieżność}
\newcommand{\lblPartIV}{Wdrażanie}
\newcommand{\lblPartV}{Python do pracy z AI}
\newcommand{\lblPartAppendices}{Dodatki}

% ---------- Tytuły ramek ----------
\newcommand{\lblNote}{Uwaga}
\newcommand{\lblWarning}{Ostrożnie}
\newcommand{\lblTrap}{Pułapka}
\newcommand{\lblCsharp}{Po stronie C\#}
\newcommand{\lblVersion}{Zależne od wersji}
\newcommand{\lblVerify}{Uruchom, zanim zaufasz}
\newcommand{\lblExercise}{Ćwiczenie}
\newcommand{\lblProject}{Etap trace-assert}
\newcommand{\lblMeasured}{Zmierzone}

% ---------- Stopka ćwiczenia ----------
\newcommand{\lblExerciseStarter}{Plik startowy}
\newcommand{\lblExerciseTest}{Test}
\newcommand{\lblExerciseRun}{Uruchom}

% ---------- Znaczniki stanu builda ----------
\newcommand{\lblNotWritten}{JESZCZE NIENAPISANE}
\newcommand{\lblNotRendered}{NIEWYRENDEROWANE --- uruchom \texttt{make diagrams}}
\newcommand{\lblSourceMissing}{Brak źródła Mermaid}
\newcommand{\lblTranscriptMissing}{TRANSKRYPT NIEWYGENEROWANY --- uruchom \texttt{make numbers}}
\newcommand{\lblValueMissing}{brak wartości}
\newcommand{\lblDiagramManifest}{Manifest diagramów}
\newcommand{\lblExerciseManifest}{Manifest ćwiczeń}

% ---------- Żywe paginy frontmatter ----------
\newcommand{\lblHowToUse}{Jak korzystać z tej książki}
\newcommand{\lblIntroduction}{Wprowadzenie}

% Myślnik interpunkcyjny. Polski składa półpauzę ze spacją po obu stronach,
% angielski pauzę bez spacji. Autorzy piszą \dash i nigdy nie wstawiają
% myślnika wprost.
\newcommand{\dash}{--}

% ---------- Odsyłacze w indeksie ----------
\newcommand{\lblIndexSee}{zob.}
\newcommand{\lblIndexSeeAlso}{zob.\ też}

% Data ostatniego sprawdzenia przypiętych wersji względem PyPI.
\newcommand{\pinnedon}{wrzesień 2026}

% ---------- Metadane PDF ----------
\newcommand{\lblPdfSubject}{Czytanie, pisanie i wdrażanie Pythona, kiedy C\# jest twoim pierwszym językiem}
\newcommand{\lblPdfKeywords}{python, dotnet, csharp, zmiana technologii, uv, pydantic, fastapi, sqlalchemy, pytest, asyncio}
\newcommand{\lblPdfLang}{pl-PL}
